% Fate's Edge SRD --- Core Principles (Section 01)

\section{Core Principles}

\subsection{Safety Tools: Lines, Veils \& X-Card}

Before play, establish how the table handles sensitive content. Fate's Edge uses three standard tools.

\subsubsection*{X-Card}
Any player (including the GM) may draw an X on a card or simply say ``X'' when a moment makes them uncomfortable. When the X-Card is shown, the described content is removed from the fiction — no questions, no explanations, no judgments. The scene is retconned or redirected immediately.

\subsubsection*{Lines}
A Line is a hard boundary: content that will not appear in the game at all. Examples: ``No sexual violence,'' ``No harm to children,'' ``No body horror.'' Before the campaign or a session, players name their Lines. The GM and group honor them completely.

\subsubsection*{Veils}
A Veil is a soft boundary: content that can happen off screen or fade to black. The event is acknowledged but not described in detail. Example: ``A romantic scene can happen, but fade to black before it becomes explicit.'' Veils allow the story to acknowledge mature themes without forcing players to roleplay through them.

\subsection{What Is Fate's Edge?}

Fate's Edge is a narrative-first tabletop roleplaying system where every action carries weight, every choice has consequence, and even failure feeds the story. Dice are not merely success meters — they are instruments of fate, weaving opportunity with risk.

\begin{quote}
\textit{``The road is a ledger. Every broken wheel is a debt, every lit lamp a witness. The only question is: what are you willing to owe?''}
\end{quote}

\subsection{The Central Question}

What are you willing to risk, and what are you willing to pay, to reshape the world around you?

\subsection{Tone of Play}

\begin{itemize}
  \item \textbf{Cinematic:} Scenes breathe; tense exchanges resolve in moments.
  \item \textbf{Consequential:} A whispered secret can topple a kingdom; a missed roll can spare a life.
  \item \textbf{Collaborative:} The GM sets stakes; the players drive the fiction.
\end{itemize}

\subsection{Key Concepts}

\subsubsection{Narrative Time}
Time is measured by story weight, not by clocks:

\begin{itemize}
  \item \textbf{A Moment:} a heartbeat, a strike, a whispered word.
  \item \textbf{Some Time:} a few minutes — a skirmish, a careful lockpick, a short negotiation.
  \item \textbf{Significant Time:} hours — travel between locations, a ritual, an interrogation.
  \item \textbf{Days:} marches across countryside, training a cadre, recovery from serious harm.
\end{itemize}

\subsubsection{Story Beats (SB)}
Every time a player rolls a 1 on a die, the GM gains a Story Beat. SB are narrative currency — spend them to introduce complications, escalate tension, shift Position, or advance clocks.

\subsubsection{Affinity}
Each culture grants an Affinity — a narrative edge or metaphysical bond that makes certain actions more reliable (e.g., Aeler stone-sense, Aelaerem hearth-luck). Affinities are described in the full cultural guides.

\subsubsection{Prestige Abilities}
High-level talents unlocked by mastering cultural or arcane arts. They are narrative milestones as much as mechanical upgrades.

\subsubsection{On-Screen vs. Off-Screen}
\begin{itemize}
  \item \textbf{On-Screen:} followers, allies, companions who stand beside you in a scene.
  \item \textbf{Off-Screen:} assets, titles, informants, safe houses — resources that act when you spend Boons or during downtime.
\end{itemize}

\subsection{Boons Overview}

Boons are narrative tokens earned primarily by failing rolls with meaningful consequences. They represent momentum, luck, or a second wind.

\textbf{Earning Boons:}
\begin{itemize}
  \item \textbf{Partial} ($0 <$ successes $<$ DV): Gain 1 Boon.
  \item \textbf{Miss} (successes = 0): Gain 2 Boons (provided the roll was meaningful and generated at least one SB).
  \item \textbf{Bond-driven action:} Once per bond per session, when you act meaningfully on a Mutual Bond and provide an Intricate Description, gain 1 Boon.
  \item \textbf{GM award:} For creative solutions, sacrifices, or outstanding roleplay.
\end{itemize}

\textbf{Spending Boons:}
\begin{itemize}
  \item Re-roll a single die after seeing the pool.
  \item Activate an on-screen Asset.
  \item Improve Position by one step before a roll.
  \item Power certain Rites or abilities as described in their text.
  \item During downtime, convert 2 Boons $\rightarrow$ 1 XP (once per session; max 2 XP per session via conversion).
\end{itemize}

\textbf{Limits:}
\begin{itemize}
  \item Maximum 5 Boons held at any time.
  \item At the end of each scene, reduce held Boons to a maximum of 2 (excess are lost).
  \item At most 2 Boons from failures per character per scene.
  \item Repetition rule: the same action with the same stakes in the same scene cannot award another Boon.
\end{itemize}

\subsection{Fail Forward: Every Roll Matters}

\begin{itemize}
  \item A \textbf{Clean Success} advances the fiction without added friction.
  \item A \textbf{Success \& Cost} achieves the goal but introduces a complication (GM spends SB).
  \item A \textbf{Partial} makes meaningful progress and gives the player a Boon.
  \item A \textbf{Miss} fails to achieve the goal, gives the player 2 Boons, and gives the GM SB to escalate.
\end{itemize}

No roll should result in ``nothing happens.'' Every outcome changes the situation.

